\section{A compact hypersurface containing a Hirzebruch surface}\label{sec:geometry}

We specify the compact example by its fan. Let $N=\ZZ^4$, and let
$v_0,\ldots,v_{10}$ be the rows of
\begin{equation}\label{eq:rays}
 \begin{pmatrix}
 1&0&0&0\\0&1&0&0\\-1&-1&0&0\\0&0&1&0\\0&0&0&1\\
 -4&-2&-1&0\\-4&-2&0&-1\\3&1&1&1\\2&1&1&1\\
 -2&-1&0&0\\-1&0&0&0
 \end{pmatrix}.
\end{equation}
The maximal cones of $\Sigma$ are generated by the following
four-tuples of rays:
\begin{equation}\label{eq:cones}
\begin{gathered}
 (0,1,3,6),\ (0,1,3,7),\ (0,1,4,5),\ (0,1,4,7),\ (0,1,5,6),\\
 (0,2,4,5),\ (0,2,4,7),\ (0,2,5,6),\ (0,2,6,7),\ (0,3,6,7),\\
 (1,3,6,10),\ (1,3,7,8),\ (1,3,8,10),\ (1,4,5,10),\\
 (1,4,7,8),\ (1,4,8,10),\ (1,5,6,10),\ (2,4,5,9),\\
 (2,4,7,9),\ (2,5,6,9),\ (2,6,7,9),\ (3,6,7,9),\\
 (3,6,9,10),\ (3,7,8,9),\ (3,8,9,10),\ (4,5,9,10),\\
 (4,7,8,9),\ (4,8,9,10),\ (5,6,9,10).
\end{gathered}
\end{equation}
Write $V_\Sigma$ for the associated toric fourfold and $D_i$ for
its invariant divisors. Let $Y$ be a general anticanonical
hypersurface, nondegenerate with respect to its Newton polytope,
with nonzero vertex coefficients. This is the resolution chamber
used in \cite{Higgs,Transition}; the construction below uses only
the displayed toric data.

\subsection{The smooth projective hypersurface}
The polytope $P=\conv(v_0,\ldots,v_8)$ is reflexive, and
\eqref{eq:cones} is a subdivision of its boundary fan. Each maximal
cone has determinant $\pm1$; the subdivision is complete.
For an explicit strictly convex support function take divisor
coefficients
\begin{equation}\label{eq:ample}
 A_i=1009(10000+h_i)+(i+1),\qquad
 h=(10,11,3,9,2,7,1,4,8,-5,-3).
\end{equation}
For each maximal cone $\sigma$, solve
$\langle m_\sigma,v_i\rangle=-A_i$ for $i\in\sigma$.
The off-cone values
$\langle m_\sigma,v_i\rangle+A_i$ are all positive, with minimum
$5047$. This proves projectivity. Repeating the same calculation
with all coefficients equal to one gives integral slopes and
nonnegative values. Thus the anticanonical divisor is basepoint-free
and the fan is a crepant refinement of the reflexive fan. Bertini
and adjunction give smoothness and $K_Y\simeq\mathcal O_Y$.
The anticanonical exact sequence, toric structure-sheaf vanishing,
and Serre duality give $H^1(\mathcal O_Y)=H^2(\mathcal O_Y)=0$.

The four vertices $v_0,v_1,v_3,v_4$ of $P$ form the standard
basis of $N$. In particular, the lattice $N''$ generated by the
lattice points on vertices and edges of $P$ is all of $N$.
The Brauer-group formula of \cite[Corollary 3.9]{BK} therefore gives
\[
 \operatorname{Br}(Y)
 \simeq\operatorname{Hom}\bigl(\Lambda^2N/(N\wedge N''),\QQ/\ZZ\bigr)=0.
\]
For this Calabi--Yau threefold the Brauer group is the torsion in
$H^3(Y;\ZZ)$. The universal coefficient theorem then implies
\begin{equation}\label{eq:torsion}
 \operatorname{Tors}H_2(Y;\ZZ)=0.
\end{equation}
This uses the full lattice, rather than a rational cohomology rank.

\subsection{Nef divisors and the contracted locus}
The divisor $D_{10}$ misses $Y$: its dual Newton face is a vertex
whose coefficient is nonzero. Define integral divisors $H_1,\ldots,H_6$
on $Y$ by the columns of
\begin{equation}\label{eq:nef-matrix}
 \begin{pmatrix}
 0&0&1&1&2&2\\
 1&1&2&4&5&5\\
 0&1&2&4&4&5\\
 0&-1&-1&-3&-3&-3\\
 0&-1&-1&-2&-2&-2\\
 0&0&1&2&2&2
 \end{pmatrix}
\end{equation}
in the ordered divisor list $(D_2,D_5,D_6,D_7,D_8,D_9)|_Y$.
Each has a nef lift to $V_\Sigma$: use the displayed coefficients,
zero on the other rays except $10$, and add coefficient zero on
$D_{10}$ for $H_1,H_2,H_3$, or coefficient one for $H_4,H_5,H_6$.
Solving the support-function equalities as above gives nonnegative
values of the support-function inequalities for every cone and every ray.
Since $D_{10}|_Y=0$, all six
$H_i$ are nef on $Y$. Toric linear equivalence identifies $H_4$
with $D_0|_Y$.

The support function of $D_0$ defines a semiample contraction.
Its positive-dimensional orbit fibers have the following minimal
cones in their closures:
\begin{equation}\label{eq:contracted-cones}
 (8),\quad(9),\quad(10),\quad(2,4),\quad(6,7).
\end{equation}
Here is a direct finite way to obtain the list. Form
$P_0=\{m:\langle m,v_i\rangle+\delta_{i0}\ge0\}$ and, for
each cone $\sigma$, impose its equalities to get the face
$F_\sigma\subset P_0$. The orbit fiber is positive-dimensional
exactly when $4-\dim F_\sigma>\dim\sigma$; taking minimal cones
in this set gives \eqref{eq:contracted-cones}.
Consequently every curve of $H_4$-degree zero on $Y$ lies in
\begin{equation}\label{eq:contracted-locus}
 E\cup F\cup C_1\cup C_2,
 \qquad E=D_9|_Y,\quad F=D_8|_Y,
\end{equation}
where $C_1,C_2$ are the intersections with the orbit closures
of $67,24$. Their restricted Newton polytopes are the primitive
segments joining $(0,1,-1,-1)$ to $(-1,2,-1,1)$ and
$(-1,2,1,-1)$, respectively. Each binomial cuts the dense torus
in a translate of $\CC^*$. Basepoint-freeness excludes boundary
components for a general equation, and nondegeneracy makes the
closures smooth; hence both curves are $\PP^1$. They miss
both surfaces: there are no cones containing $678$ or $248$,
and the Newton faces at $679,249$ are singletons. There is no
cone containing $2467$, so the two curves are disjoint.

\begin{lemma}\label{lem:surfaces}
The surfaces are $E\simeq\Bl_{p_1,p_2}\PP^2$ and $F\simeq\FF_1$.
Their intersection is the negative section $s$ on $F$, and the
strict transform of the line through $p_1,p_2$ on $E$.
\end{lemma}
\begin{proof}
The anticanonical face at ray $9$ is the primitive edge from
$(0,1,-1,-1)$ to $(1,-1,-1,-1)$. Its binomial equation fixes a
nonzero torus parameter. Rays $2,7$ are excluded on $E$ because
their intersections have singleton Newton faces; ray $10$ is absent.
The remaining rays $3,4,5,6,8$ project to
\[
 (1,0),\quad(0,1),\quad(-1,0),\quad(0,-1),\quad(1,1),
\]
with cones $38,84,45,56,63$. This is the fan of the plane blown
up at two distinct torus fixed points.

For $F$, the star of ray $8$ away from $D_{10}$ has cones
$137,147,379,479$. Quotient by the fiber ray $7$ gives the base
rays $(1,0),(0,1),(-1,-1),(-1,0)$, the fan of $\FF_1$.
The Newton vertex $(1,-1,-1,-1)$ has exponents
$(0,0,0,1,0,0)$ on rays $1,3,4,7,9,10$; all other vertices
have fiber exponent zero. The equation is monic linear in that
fiber coordinate and its zero locus is a graph over $\FF_1$.
The common divisor in these two surface fans is the stated curve.
\end{proof}

Let $e_1,e_2$ denote the exceptional curves on $E$, and $f$ a
ruling fiber of $F$. Restriction of \eqref{eq:nef-matrix} to the
surface fans and to the two orbit curves gives
\begin{equation}\label{eq:degrees}
\begin{array}{c|rrrrrr}
 &H_1&H_2&H_3&H_4&H_5&H_6\\\hline
 s&0&1&0&0&0&0\\
 f&0&0&1&0&0&0\\
 e_1&1&0&0&0&1&1\\
 e_2&0&0&0&0&0&1\\
 C_1&0&0&0&0&1&0\\
 C_2&1&0&0&0&0&0.
\end{array}
\end{equation}
These pairings can be computed with the usual toric divisor
relations: the linear relations are
$\sum_i\langle m,v_i\rangle D_i=0$, and a divisor product vanishes
when its rays span no cone. The monic and binomial equations above
identify the actual surfaces on which the restrictions are taken.

\begin{proposition}\label{prop:support}
The class $\alpha=3s+4f$ is primitive in $H_2(Y;\ZZ)$.
Every effective cycle of class $m\alpha$, $m\ge1$, is supported on
$F$. The same holds for every effective additive summand of $\alpha$.
\end{proposition}
\begin{proof}
By \eqref{eq:degrees}, $(H_3-H_2)\cdot\alpha=4-3=1$,
which proves primitivity. A cycle in class $m\alpha$ has zero
$H_1,H_4,H_5,H_6$ degrees. Nefness forces the same for each
component. The $H_4$ contraction confines it to
\eqref{eq:contracted-locus}; $H_1,H_5$ exclude $C_1,C_2$.
The effective cone of $E$ is generated by $e_1,e_2,s$.
The $H_1,H_6$ degrees exclude any $e_1$ or $e_2$ coefficient,
so a remaining component in $E$ is supported on $s=E\cap F$.
All components therefore lie in $F$.

Every effective summand of $\alpha$ has the same four vanishing
degrees, so the argument applies to it as well. On $F$, the
$H_2,H_3$ pairings recover the coefficients of $s,f$. Hence a
decomposition of $\alpha$ into effective classes corresponds to
an actual effective decomposition of $3s+4f$ on $F$.
\end{proof}
